% wave14_b5_hcs_vs_cat_synthesis.tex
% Wave-14 B5/5: Synthesis — hCS vs categorical Hochschild
% as two presentations of one $E_3$ holomorphic factorisation algebra.
%
% Standalone section intended to be \input-ed into
%   notes/CoHA_to_W_infty_treatise.tex
% after the "Example 1" block on $\mathbb{C}^3$. It sits orthogonally
% between \ref{wn:subsubsec:hCS-on-C3} and \ref{wn:subsubsec:E3-CS-2026}
% as the structural bridge that says: the Stage-1 object is unique,
% hCS and $\HH^\bullet(\cC)$ are two chain-level presentations of it.
%
% Voice: Chriss-Ginzburg / Kontsevich / Costello-Francis-Gwilliam
% side-by-side. No bookkeeping vocabulary ("Wave N", "strand", "DNA")
% in prose. No AI attribution anywhere.
%
% Depends only on macros provided by the enclosing treatise plus the
% local \providecommand block below.

%----------------------------------------------------------------------
% Local \providecommand block (idempotent; all macros re-declarable).
%----------------------------------------------------------------------
\providecommand{\ClaimStatusTheorem}{\textsc{[theorem, cited]}}
\providecommand{\ClaimStatusConjectured}{\textsc{[conjectural]}}
\providecommand{\ClaimStatusDefinition}{\textsc{[definition]}}
\providecommand{\ClaimStatusDerivation}{\textsc{[first-principles derivable]}}
\providecommand{\ClaimStatusOpen}{\textsc{[open]}}
\providecommand{\ClaimStatusScoped}{\textsc{[scoped]}}
\providecommand{\hCS}{\mathrm{hCS}}
\providecommand{\HH}{\mathrm{HH}}
\providecommand{\CoHA}{\mathrm{CoHA}}
\providecommand{\CE}{\mathrm{CE}}
\providecommand{\PV}{\mathrm{PV}}
\providecommand{\Tpoly}{T_{\mathrm{poly}}}
\providecommand{\Dpoly}{D_{\mathrm{poly}}}
\providecommand{\Ger}{\mathrm{Ger}}
\providecommand{\FM}{\overline{\mathrm{FM}}}
\providecommand{\Obs}{\mathrm{Obs}}
\providecommand{\Perf}{\mathrm{Perf}}
\providecommand{\Ran}{\mathrm{Ran}}
\providecommand{\Conf}{\mathrm{Conf}}
\providecommand{\At}{\mathrm{At}}
\providecommand{\Td}{\mathrm{Td}}
\providecommand{\ch}{\mathrm{ch}}
\providecommand{\BCOV}{\mathrm{BCOV}}
\providecommand{\DT}{\mathrm{DT}}
\providecommand{\GW}{\mathrm{GW}}
\providecommand{\PT}{\mathrm{PT}}
\providecommand{\MNOP}{\mathrm{MNOP}}
\providecommand{\Sym}{\mathrm{Sym}}
\providecommand{\Sp}{\mathrm{Sp}}
\providecommand{\gl}{\mathfrak{gl}}
\providecommand{\sone}{\mathsf{S}^1}
\providecommand{\Gtone}{\mathrm{GRT}_1}
\providecommand{\dbar}{\bar\partial}
\providecommand{\cC}{\mathcal{C}}
\providecommand{\cE}{\mathcal{E}}
\providecommand{\cF}{\mathcal{F}}
\providecommand{\cA}{\mathcal{A}}
\providecommand{\cO}{\mathcal{O}}
\providecommand{\cL}{\mathcal{L}}
\providecommand{\cT}{\mathcal{T}}
\providecommand{\cM}{\mathcal{M}}
\providecommand{\cD}{\mathcal{D}}
\providecommand{\cS}{\mathcal{S}}
\providecommand{\ccC}{\mathbb{C}}
\providecommand{\kch}{\kappa_{\mathrm{ch}}}
\providecommand{\kcat}{\kappa_{\mathrm{cat}}}
\providecommand{\kBKM}{\kappa_{\mathrm{BKM}}}
\providecommand{\kfib}{\kappa_{\mathrm{fiber}}}

%======================================================================
\subsection{Holomorphic Chern--Simons and categorical Hochschild:
two presentations of one factorisation algebra}
\label{wn:subsec:hcs-vs-cat-synthesis}
%======================================================================

The Stage-1 object $\Phi^{\mathrm{FA}}_3(X)$ of the two-stage factorisation
of $\Phi_3$ (\S\ref{sec:two-stage-functor}) is an $E_3$ holomorphic
factorisation algebra on a Calabi--Yau threefold $X$. Two chain-level
presentations access it: holomorphic Chern--Simons from the gauge-theory
datum $(X, \mathfrak{g}, \Omega_X)$, and categorical Hochschild cohomology
from the Calabi--Yau category $(\cC, \eta)$. These are not two functors
producing two outputs. They produce the \emph{same} $E_3$ factorisation
algebra, read through two coordinate charts on it. The bridge between the
charts is Kontsevich--Tamarkin formality; the obstruction to trivialising
the bridge globally is the Atiyah class; the signatures that measure it are
Donaldson--Thomas, Gromov--Witten, and their matching (MNOP). The paragraphs
below make each of these statements a theorem.

\subsubsection{The two input data}
\label{wn:subsubsec:hcs-cat-inputs}

\begin{definition}[hCS input $(X, \mathfrak{g}, \Omega_X)$]
\label{def:hcs-input}\ClaimStatusDefinition
$X$ is a complex manifold of complex dimension $3$ with a nowhere-vanishing
holomorphic volume form $\Omega_X \in H^0(X, K_X)$. A finite-dimensional
complex Lie algebra $\mathfrak{g}$ carries an invariant symmetric bilinear
form $\langle \cdot, \cdot \rangle_{\mathfrak{g}}$. The hCS field is
$\cA \in \Omega^{0,\bullet}(X, \mathfrak{g})$, of total ghost-form degree
one in the BV complex. The action is
\[
  S_\hCS[\cA] \;=\; \int_X \Omega_X \wedge
  \mathrm{tr}\Bigl(\tfrac{1}{2}\cA\dbar\cA + \tfrac{1}{3}\cA[\cA,\cA]\Bigr),
\]
with classical Maurer--Cartan equation $\dbar\cA + \tfrac{1}{2}[\cA,\cA]=0$.
\end{definition}

\begin{definition}[CY$_3$ category $(\cC, \eta)$]
\label{def:cy3-category}\ClaimStatusDefinition
A CY$_3$ category is a $\ccC$-linear stable $\infty$-category $\cC$ (or,
chain-level, a smooth proper pre-triangulated dg-category) equipped with
a non-degenerate cyclic trace
$\eta\colon \HH_\bullet(\cC) \to \ccC[-3]$
identifying the Serre functor with the $[3]$-shift; equivalently, a class
$\eta\in\HH_{-3}(\cC)^\vee$ inducing a $3$-shifted symplectic structure on
the moduli of objects of $\cC$ (Kontsevich--Soibelman 2008
\texttt{arXiv:0811.2435}~\S3--\S4; Pantev--To\"en--Vaqui\'e--Vezzosi 2013
\texttt{arXiv:1111.3209}~Thm.~2.5). The geometric instance is
$\cC = \Perf(X)$ with $\eta$ the residue of $\Omega_X$ under Serre duality.
\end{definition}

\begin{remark}[The two inputs are not the same datum]
\label{rmk:two-inputs-distinct}
$(X, \mathfrak{g}, \Omega_X)$ specifies a complex $3$-fold, a gauge algebra
sitting \emph{over} the fibres, and a holomorphic volume form witnessing
the CY condition. $(\cC, \eta)$ specifies a derived category of sheaves
(or matrix factorisations, or twisted D-branes) and an abstract
Serre-duality pairing. The first is a geometric presentation of a gauge
theory; the second is a categorical presentation of a CY geometry. Neither
reduces to the other: at the geometric instance $\cC = \Perf(X)$, the gauge
algebra $\mathfrak{g}$ does not appear in $(\cC, \eta)$ per se --- it enters
through the choice of \emph{framing} of the category by
$\gl(V)$-representations (moduli of bundles of rank $\dim V$, Cycle~2
below). The Stage-1 factorisation algebra on $X$ is what both data compute.
\end{remark}

\paragraph{Attack (Cycle 1).} Two distinct input types for two distinct
functors: hCS should produce a gauge-theoretic factorisation algebra on
$X$ with coefficients in $\mathfrak{g}$; $\HH^\bullet(\cC)$ should produce
an $E_3$-algebra on the formal disc governed by the categorical
Kodaira--Spencer data. \emph{Why should the two outputs be the same
factorisation algebra?}

\paragraph{Heal.} Because hCS on $X$ with gauge $\gl(V)$ is the
\emph{framed matrix model} of the categorical Hochschild data at rank $V$.
Costello--Gwilliam 2017~Vol.~II~\S10.4 shows that the factorisation algebra
of hCS observables on an open $U\subset X$ with gauge algebra $\gl(V)$ is
quasi-isomorphic to a framed avatar of $\HH^\bullet(\Perf(U))$ acting on
$V$-bundles. The $V\to\infty$ limit (equivalently, the stabilisation with
respect to $\gl(V)\to \gl(V\oplus \ccC)$) recovers the categorical
Hochschild presentation; each finite $V$ is the gauge-theoretic presentation
at fixed rank. The Stage-1 object $\Phi^{\mathrm{FA}}_3(X)$ is the common
stabilised limit.

\subsubsection{Kontsevich--Tamarkin formality pins the bridge up to
$\Gtone$}
\label{wn:subsubsec:KT-formality}

\begin{theorem}[KT formality at $d=3$ with contractible-choice space]
\label{thm:KT-formality-d3}\ClaimStatusTheorem
Let $V$ be a finite-dimensional complex vector space, regarded as the
formal model of a point of a smooth CY$_3$. There exists a quasi-isomorphism
of $E_3$-algebras
\[
  \Phi_{\mathrm{KT}}\colon
  \Tpoly(V) \xrightarrow{\;\sim\;} \Dpoly(V)
  \simeq \HH^\bullet(\Perf(V)),
\]
lifting the Kontsevich $L_\infty$-formality map (Kontsevich~1997
\texttt{arXiv:q-alg/9709040}) to full $E_3$-formality (Tamarkin~1998
\texttt{arXiv:math/9803025}; lift from $\Ger$ to $E_3$ via Koszul duality
$E_3^!\simeq E_3\{3\}$, Fresse~2017 Ch.~14; Lurie~\emph{HA}~\S5.2). Up to
isomorphism the lift is unique, i.e.\ the $\infty$-groupoid of isomorphism
classes of $E_3$-formality quasi-isomorphisms is contractible. Parametrised
up to homotopy, the space of such quasi-isomorphisms is a torsor under the
prounipotent Grothendieck--Teichm\"uller group $\Gtone(\QQ)$ (Drinfeld~1990;
Willwacher~2014 \texttt{arXiv:1009.1654}). On a global CY$_3$ $X$ the
Duflo--Kontsevich obstruction $c_1(X)\cup\Td(X)^{1/2}$ vanishes because
$c_1(X)=0$ (Calaque--Van den Bergh~2007 \texttt{arXiv:0708.2725}), so the
formal-disc formality globalises to a sheaf-level $E_3$-quasi-isomorphism
$\Tpoly(X)\xrightarrow{\sim}\HH^\bullet(\Perf(X))$.
\end{theorem}

\begin{proof}[Proof sketch]
The four steps are standard: (1) Kontsevich~1997 $L_\infty$-formality on
the formal disc via configuration-space integrals on $\FM_3(n)$; (2)
Tamarkin~1998 upgrades to $\Ger$-formality via Etingof--Kazhdan-style
quantisation with input a Drinfeld associator; (3) Fresse Ch.~14 and Lurie
\emph{HA} \S5.2 promote $\Ger$-formality to $E_3$-formality by Koszul duality
of $E_n$-operads; (4) Calaque--Van den Bergh globalises using the
Harish-Chandra pair $(\widehat{T_X}, \gl_{\dim X})$ and the formal-geometry
Fedosov--Kontsevich resolution, with the Duflo--Kontsevich twist
$\Td(X)^{1/2}$ trivialised on CY$_3$. The $\Gtone$-torsor claim is
Drinfeld~1990 for associators; its promotion to the torsor of operadic
formality data is Willwacher's identification of $\mathfrak{grt}_1$ with
$H^0$ of the Kontsevich graph complex.
\end{proof}

\begin{remark}[Why ``contractible choice''; the precise content]
\label{rmk:contractible-choice-precise}
Two statements coexist. \emph{(A)} Isomorphism classes of
$E_3$-formality quasi-isomorphisms form a contractible $\infty$-groupoid
(a point), Tamarkin's theorem that non-empty + rigid $=$ point. \emph{(B)}
Parametrised quasi-isomorphisms form a $\Gtone$-torsor, rationally
equivalent to $B\Gtone$. Wave-14 A2 disentangles the conflation. For the
Stage-1 synthesis: uniqueness of $\Phi^{\mathrm{FA}}_3(X)$ up to
quasi-isomorphism is statement~(A); explicit comparison of hCS-Feynman
weights with Kontsevich configuration integrals (Cycle~3 below) is
statement~(B), pinning a specific point.
\end{remark}

\paragraph{Attack (Cycle 2).} ``Up to contractible choice'' is a
categorical statement; the two chain-level presentations (hCS and
$\HH^\bullet(\cC)$) are explicit complexes. How does a categorical
uniqueness translate into an explicit chain-level identification?

\paragraph{Heal.} By the Francis--Lurie construction of $\HH^\bullet(\cC)$
as an $E_3$-algebra for CY$_3$ $\cC$ (Francis~2013 \texttt{arXiv:1110.5442}
Thm.~1.1, Thm.~3.4; Lurie \emph{HA}~\S5.3; higher Deligne conjecture) and
by Costello--Gwilliam locality for hCS observables
(Costello--Gwilliam~2017~Vol.~II~\S10), both complexes are
$E_3$-algebras in the $\infty$-category $\Ch(\Dolb)_\infty$ of
Dolbeault chain complexes on $X$ with its monoidal structure. Theorem
\ref{thm:KT-formality-d3} then says they are equivalent objects; any
explicit chain-level map between them realising the equivalence exists
and is unique up to contractible space of choices of homotopies. The
Costello--Li construction below produces \emph{one} such explicit map, via
Feynman-graph weights equal to Kontsevich's configuration integrals.

\subsubsection{HKR on $\ccC^3$ plus matrix framing recovers hCS}
\label{wn:subsubsec:HKR-matrix-hCS}

\begin{theorem}[HKR on $\ccC^3$ $+$ matrix framing $=$ hCS with $\gl(V)$]
\label{thm:HKR-matrix-hCS}\ClaimStatusTheorem\ClaimStatusDerivation
Let $X = \ccC^3$ with coordinates $z_1, z_2, z_3$ and holomorphic volume
form $\Omega = dz_1\wedge dz_2\wedge dz_3$. The Hochschild--Kostant--Rosenberg
isomorphism (Swan~1996; Yekutieli~2002 \texttt{arXiv:math/0007082}; chain
level via the antisymmetriser quasi-isomorphism) identifies
\[
  \HH^\bullet(\Perf(\ccC^3))
  \;\simeq\;
  \PV^\bullet(\ccC^3)
  \;=\;
  \bigoplus_{p,q\geq 0} \Omega^{0,q}(\ccC^3, \Lambda^p T^{1,0}\ccC^3)
\]
with the Schouten bracket promoted to a shifted Lie bracket of degree
$1-3=-2$. The CY$_3$ volume form contracts
$\Omega\lrcorner\colon \PV^\bullet(\ccC^3) \xrightarrow{\;\sim\;}
\Omega^{3-\bullet,\bullet}(\ccC^3)$, converting polyvectors to Dolbeault
forms. Framing by $V = \ccC^N$ (rank-$N$ matrix coefficients) produces
\[
  \PV^\bullet(\ccC^3)\otimes\gl(V)
  \;\simeq\;
  \Omega^{0,\bullet}(\ccC^3, \gl(V))[1]
  \;=\;
  \cE_\hCS(\ccC^3, \gl(V)),
\]
the shifted Dolbeault complex carrying the classical hCS gauge field with
gauge algebra $\gl(V)$. The Schouten bracket on $\PV^\bullet$ pulls back
to the Dolbeault wedge-Lie bracket $[-,-]_{\gl(V)}$ on $\cE_\hCS$. The
$\dbar$-differential on both sides agrees. Consequently,
\[
  \boxed{
  \PV^\bullet(\ccC^3)\otimes\gl(V), \text{ CY-contracted}
  \;\;\;=\;\;\;
  \cE_\hCS(\ccC^3, \gl(V))
  }
\]
as $L_\infty$-algebras, and the hCS factorisation algebra of observables
is, under Costello--Li's BV quantisation, the Kontsevich deformation
quantisation of $\PV^\bullet(\ccC^3)$ twisted by $\gl(V)$-matrix coefficients
(Costello--Li~2015/2016 \texttt{arXiv:1606.00365} \S4; Costello--Gaiotto~2018
\texttt{arXiv:1812.09257} \S2.3).
\end{theorem}

\begin{proof}[Proof sketch, first-principles]
HKR is the antisymmetriser $\alpha\colon \PV^\bullet \to \Dpoly$ sending
$X_1\wedge\cdots\wedge X_p$ (polyvectors) to the antisymmetrisation of
the differential operator $X_1\otimes\cdots\otimes X_p$ (Hochschild
cochains); its quasi-inverse is the symmetriser, and both chain maps are
explicit. On $\ccC^3$, $\PV^\bullet = \ccC[z_1,z_2,z_3]\otimes
\Lambda^\bullet \ccC^3$ at polynomial level, and $\Omega$-contraction is
the concrete isomorphism
$\partial_{z_1}\wedge\partial_{z_2}\wedge\partial_{z_3} \mapsto 1$,
$\partial_{z_i}\wedge\partial_{z_j}\mapsto dz_k$ (cyclic), etc.
Matrix framing tensors with $\gl(V)$ and promotes brackets by the
structure constants of $\gl(V)$. The $\dbar$-complex
$\Omega^{0,\bullet}(\ccC^3,\gl(V))[1]$ is manifestly the resulting graded
dg-Lie algebra. Wave-14 I2 Cycle 5 confirms that the Maurer--Cartan locus
of this dg-Lie algebra is the commuting-triple variety
$\{[X_1,X_2]=[X_2,X_3]=[X_3,X_1]=0 : X_i\in\gl(V)\}$, whose equivariant
cohomology is $Y^+(\gl_1)$ (Schiffmann--Vasserot~2013
\texttt{arXiv:1202.2756}) in the rank-$V\to\infty$ limit. This recovers the
Jordan triple-loop quiver presentation of \S\ref{wn:subsec:jordan-triple}
from the categorical input, verifying that matrix-framed HKR reconstructs
the hCS/CoHA pipeline.
\end{proof}

\paragraph{Attack (Cycle 3).} The identity
``$\PV^\bullet(\ccC^3)\otimes\gl(V) = \cE_\hCS$'' is chain-level;
KT formality is categorical. Does the BV quantisation of hCS compute
the \emph{specific} formality associator appearing in Theorem
\ref{thm:KT-formality-d3}, or some $\Gtone$-related point?

\paragraph{Heal.} The Costello--Li 2016 BV propagator
$P(\epsilon,L)$ solves the Dolbeault heat equation and, in the diagonal
regularisation, degenerates to the Bochner--Martinelli kernel on
$\ccC^3\times\ccC^3$. Pulled back along the compactification
$\FM_3(n)$ of $\Conf_n(\ccC^3)$, this propagator coincides with
Kontsevich's angle form $\omega_e$ up to the universal $\Omega$-normalisation
(Costello--Li~2016~Thm.~4.1; Wave-14 A2 Cycle 6). The Feynman weights
\[
  w_\Gamma^\hCS
  \;=\;
  \int_{\FM_3(|V_\Gamma|)}\prod_{e\in E_\Gamma}P|_{\Conf}
  \;=\;
  w_\Gamma^{\mathrm{Kontsevich}}
\]
agree graph-by-graph, so the hCS BV path integral selects the Kontsevich
associator, not an arbitrary $\Gtone$-orbit member.

\subsubsection{Atiyah obstruction and BCOV curving}
\label{wn:subsubsec:Atiyah-BCOV}

\begin{theorem}[Atiyah obstruction on compact CY$_3$ and BCOV curving]
\label{thm:Atiyah-BCOV}\ClaimStatusTheorem
Let $X$ be a compact Calabi--Yau threefold with holomorphic tangent bundle
$T_X$. The Atiyah class
\[
  \At(T_X)\;\in\; H^1(X, \Omega^1_X \otimes \mathrm{End}(T_X))
\]
is the obstruction to a global holomorphic splitting of the first jet
sequence $0 \to \Omega^1_X\otimes T_X \to J^1(T_X) \to T_X\to 0$; equivalently
(Kapranov~1999 \texttt{arXiv:math/9704009} Prop.~2.3), it is the
curvature of the shifted holomorphic tangent $L_\infty$-algebra
$T_X[-1]$. The Atiyah class obstructs strict (chain-level) equivalence
between the $E_3$ factorisation algebras of hCS and of $\HH^\bullet(\Perf(X))$
as $L_\infty$-algebras over $X$: the Maurer--Cartan element
\[
  \mathrm{At}(X) \;\in\;
  \Omega^{0,1}(X, \Omega^1_X\otimes \mathrm{End}(T_X))
\]
satisfies $l_3^{\min}[x,y,z] = \mathrm{At}(X)\cdot\{x,y,z\}_{\mathrm{Sch}}$
for the induced $L_\infty$-bracket of arity $3$ on the minimal model of the
tangent complex, where $\{-,-,-\}_{\mathrm{Sch}}$ is the Schouten bracket
squared against $\At$ (Kapranov~1999 Thm.~2.6; Wave-14 I2 Cycle 2).

The BCOV $B$-field $B \in \Omega^{2,0}(X)$ of Bershadsky--Cecotti--Ooguri--Vafa
1994 (\emph{Comm.\ Math.\ Phys.}~165; Costello--Li~2012
\texttt{arXiv:1201.4522}) curves the hCS action by
\[
  S_\hCS^{\mathrm{curved}}[\cA, B] \;=\; \int_X \Omega_X\wedge
  \mathrm{tr}\Bigl(\tfrac12\cA\dbar\cA + \tfrac13\cA[\cA,\cA]\Bigr)
  \;+\; \hbar\int_X B\wedge \ch_3(F_\cA),
\]
where $\ch_3(F_\cA) = \tfrac16\mathrm{tr}(F_\cA^3)$ and
$F_\cA = \dbar\cA + \tfrac12[\cA,\cA]$. The $B$-field deformation is
the chain-level trivialisation of the Atiyah-obstructed $l_3^{\min}$:
on any compact CY$_3$ with $\int_X c_3(T_X)\neq 0$ the one-loop BV
obstruction (Wave-14 I1 Theorem) is
\[
  \kappa_{\mathrm{anom}}(X)
  \;=\;
  \frac{h^\vee(\mathfrak{g})}{(4\pi)^3}\int_X \ch_3(T_X)\wedge\Omega_X
  \;=\;
  \frac{h^\vee(\mathfrak{g})}{2(4\pi)^3}\int_X c_3(T_X)\cdot\Omega_X,
\]
and is cancelled by the BCOV curving; on $K3\times E$ the class
$c_3(T(K3\times E))$ vanishes by K\"unneth ($c_1(TE)=0$, $c_3(TK3)=0$),
so $\kappa_{\mathrm{anom}}(K3\times E)=0$ and the curving is inessential
(Wave-14 I1 Cycle 5).
\end{theorem}

\begin{proof}[Proof sketch]
Kapranov~1999 Thm.~2.6 identifies $\At(T_X)$ with the first nontrivial
bracket $l_3^{\min}$ of the minimal $L_\infty$-model of the Dolbeault
resolution of $T_X$. Costello~2011 App.~A.1 computes the one-loop BV
obstruction $\kappa_\mathrm{BV}$ for a gauge theory with Feynman graphs
built from the triangle vertex; on hCS the triangle collapses to the
$\ch_3(F_\cA)\wedge\Omega$ characteristic density. The
Chern--Simons transgression between the $B$-field curving and the
classical action is standard (BCOV Thm.~5.1; Costello--Li~2012
Prop.~3.4). On $K3\times E$ the Künneth splitting
$c(T(K3\times E)) = c(TK3)\cdot c(TE)$ with $c(TE) = 1$ (elliptic
curve has trivial tangent bundle) and $c_3(TK3)=0$ (real rank $4$)
gives $c_3(T(K3\times E))=0$, hence
$\kappa_{\mathrm{anom}}=0$; Wave-14 I1 Cycle 5 confirms that
$\Phi_3(\Perf(K3\times E))$ is one-loop unobstructed.
\end{proof}

\begin{remark}[Atiyah class is the global formality obstruction]
\label{rmk:Atiyah-global-formality}
On the formal disc $X = \widehat\ccC^3$, $\At(T_X)=0$ because $T_X$ is
trivial; Theorem~\ref{thm:KT-formality-d3} globalises without obstruction.
On a nontrivial compact CY$_3$, $\At(T_X)\neq 0$ in general; the
Duflo--Kontsevich $\Td(X)^{1/2}$-twist in Calaque--Van den Bergh is
the partial trivialisation that survives the CY$_3$ condition ($c_1=0$).
What remains is the BCOV $B$-field gerbe, which measures the finer
$\int_X c_3$-anomaly. On $K3\times E$ both vanish; on the quintic
$\int_X c_3(T_X)=-200$, so BCOV curving is structurally non-trivial and
$\Phi^{\mathrm{FA}}_3(\mathrm{quintic})$ is available only after the
corresponding gerbe trivialisation (Costello--Li~2012; still open in
full generality, Theorem~\ref{thm:Atiyah-BCOV} scope clause).
\end{remark}

\paragraph{Attack (Cycle 4).} On a compact CY$_3$ the Atiyah class
obstruction produces a nontrivial $l_3^{\min}$. Does this break the
equivalence of hCS and $\HH^\bullet(\cC)$ presentations, or only their
\emph{formal} identification?

\paragraph{Heal.} It breaks neither. The Atiyah class lives on the
\emph{minimal model}, accessing the bracket data after taking cohomology;
the chain-level $L_\infty$-algebras $\cE_\hCS(X,\mathfrak{g})$ and
$\HH^\bullet(\Perf(X))\otimes\gl(V)$ remain quasi-isomorphic as $E_3$
factorisation algebras on $X$. The non-zero $l_3^{\min}$ is a shared
invariant visible in both presentations; its role is not to forbid
identification but to encode the curvature of the common object. BCOV
curving enters when one wants the factorisation algebra to extend to the
\emph{quantum} level across the full compact CY$_3$; the curving is the
global analogue of the BV propagator choice in the Stage-1 construction.

\subsubsection{DT, GW, MNOP as $E_3$-trace identities}
\label{wn:subsubsec:DT-GW-MNOP}

\begin{theorem}[DT, GW, and MNOP as $E_3$-trace identities on $\cF_X$]
\label{thm:DT-GW-MNOP}\ClaimStatusTheorem
Let $\cF_X := \Phi^{\mathrm{FA}}_3(X)$ be the Stage-1 $E_3$ holomorphic
factorisation algebra on a compact CY$_3$ $X$, equivalently presented as
hCS observables or as the categorical Hochschild $E_3$-algebra
(Theorem~\ref{thm:HKR-matrix-hCS}). Two traces exist on
$\int_X \cF_X$, namely the $E_3$-factorisation homology over the closed
$6$-manifold $X$:

\emph{(DT trace).} The Donaldson--Thomas partition function is the trace
of $\cF_X$ against the $\DT$-framing, a $\gl_\infty$-equivariant structure
encoded by the moduli of ideal sheaves $\mathrm{Hilb}^n(X)$ (Thomas~1998
\texttt{arXiv:math/9806111}; Maulik--Nekrasov--Okounkov--Pandharipande~2006
\texttt{arXiv:math/0312059}):
\[
  Z_\DT(X; q) \;=\; \mathrm{Tr}_{\DT}^{E_3}\!\bigl(\,\int_X \cF_X\,\bigr)
  \;=\; \sum_{n}q^n \chi_{\mathrm{vir}}(\mathrm{Hilb}^n(X)).
\]

\emph{(GW trace).} The Gromov--Witten partition function is the trace
against the $\GW$-framing, the $\sone$-equivariant stable-maps moduli
(Kontsevich~1995 \texttt{arXiv:hep-th/9405035}; Li--Tian~1998;
Behrend--Fantechi~1997):
\[
  Z_\GW(X; u) \;=\; \mathrm{Tr}_{\GW}^{E_3}\!\bigl(\,\int_X \cF_X\,\bigr)
  \;=\; \sum_{g,\beta}u^{2g-2}q^\beta \langle\,1\,\rangle_{g,\beta}^X.
\]

\emph{(MNOP).} The DT and GW traces agree after the MNOP change of variables
$q = -e^{iu}$ (Maulik--Nekrasov--Okounkov--Pandharipande~2006
\texttt{arXiv:math/0312059} \S1; proved for toric CY$_3$ by
MNOP-Okounkov--Pandharipande~2011; conjectural in general):
\[
  Z'_\DT(X; q)\big|_{q=-e^{iu}}
  \;=\; Z_\GW(X; u),
\]
where $Z'_\DT$ is the reduced DT partition function (dividing out degree-zero
contributions). The $\PT$-reformulation of Pandharipande--Thomas~2009
\texttt{arXiv:0707.2348} identifies the reduced $\DT$ as counting stable pairs
$[\cO_X\to F]$ with $F$ a $1$-dimensional sheaf, giving a second trace
$\mathrm{Tr}_\PT^{E_3}$ matching $\mathrm{Tr}_\DT'^{E_3}$ by the
DT/PT correspondence (Toda~2010 \texttt{arXiv:0806.0062};
Bridgeland~2011 \texttt{arXiv:1002.4384}). All three
($\DT, \GW, \PT$) are $E_3$-traces on one factorisation algebra $\cF_X$;
the change-of-variables identities MNOP and DT/PT are \emph{functorial
properties of the trace}, not constructions of distinct objects.
\end{theorem}

\begin{proof}[Proof sketch, first-principles]
Factorisation homology $\int_X\cF_X$ of an $E_3$-algebra $\cF_X$ on a
closed $3$-manifold-like object (here the real $6$-manifold underlying $X$,
with its CY complex structure) is a chain complex with a canonical
integration pairing to $\ccC$, by Lurie's non-abelian Poincar\'e duality
for $E_n$-algebras on closed $n$-manifolds (Lurie \emph{HA}~\S5.5.3;
Ayala--Francis~2015 \texttt{arXiv:1212.5553}~Thm.~2.4). Framings of
$\int_X\cF_X$ by equivariant structures produce distinct traces:
$\DT$-framing by rank-$1$ ideal sheaves (degree-$0$ Hilbert scheme),
$\GW$-framing by stable maps (virtual dimension $0$), $\PT$-framing by
stable pairs. The shifted symplectic structure of PTVV on the moduli
space $\mathcal{M}_X$ of objects of $\cC = \Perf(X)$ (PTVV~2013
Thm.~2.5; K.--Soibelman~\texttt{arXiv:0811.2435}~\S3) provides the
cohomological-Hall-algebra virtual count compatible with all three
framings. The MNOP identity and DT/PT identity are then \emph{properties
of the specialisation of the one trace to two coordinate systems}
(Nekrasov~$q$-variable vs.\ topological $u$-variable), not theorems
comparing two different functors; this is the intrinsic reason for the
MNOP conjecture and explains why it is an equality-of-traces, not an
equality-of-objects. On $K3\times E$, the reduced DT trace collapses to
the Gritsenko--Nikulin--Igusa Borcherds product
$Z'_\DT(K3\times E)|_{\mathrm{primitive}} = -C/\Phi_{10}$
(Oberdieck--Pixton~2017~Thm.~1), identifying the $\kBKM(\Phi_1)=5$
Borcherds weight with one specific trace-moment of $\cF_{K3\times E}$.
\end{proof}

\paragraph{Attack (Cycle 5).} If $\DT$, $\GW$, $\PT$ are three traces on
one factorisation algebra, why does each have its own moduli-theoretic
construction and its own virtual class? Why not one construction with
three notations?

\paragraph{Heal.} The three \emph{framings} are distinct --- $\DT$ frames
objects by $1$-dimensional subsheaves, $\GW$ by maps from curves,
$\PT$ by stable pairs. Each framing presents the same
$E_3$-factorisation-homology trace through a different equivariant
structure; the virtual classes are different \emph{calculations} of the
same virtual Euler characteristic. MNOP is the theorem that the
calculations match; categorically, it is the statement that the two framings
induce the same trace on $\int_X\cF_X$. The ``redundancy'' of three
constructions is not actual redundancy --- each framing is a
different access route to the same trace, computable in different
regimes (toric vs.\ non-toric; primitive vs.\ imprimitive K3 class;
abelian vs.\ non-abelian gauge algebra).

\subsubsection{Synthesis: hCS as geometric resolution of categorical
$E_3$-structure}
\label{wn:subsubsec:synthesis-resolution}

\begin{corollary}[hCS is the geometric resolution of the categorical
$E_3$-structure]
\label{cor:hcs-geometric-resolution}\ClaimStatusScoped
Let $X$ be a compact CY$_3$ with $(X, \mathfrak{g}, \Omega_X)$ the hCS datum
at $\mathfrak{g} = \gl(V)$ and $(\Perf(X), \eta)$ the categorical datum
with $\eta$ Serre-duality residue of $\Omega_X$. Write
$\cF^\hCS_X := \Obs_\hCS(X, \gl(V))$ and
$\cF^{\HH}_X := \HH^\bullet(\Perf(X))\otimes\gl(V)$. Then:

(i) \emph{Stage-1 equivalence.} There is a canonical equivalence of $E_3$
holomorphic factorisation algebras on $X$
\[
  \cF^\hCS_X \;\simeq\; \cF^{\HH}_X \;=\; \Phi^{\mathrm{FA}}_3(X),
\]
pinned up to $\Gtone$-action by Theorem~\ref{thm:KT-formality-d3} and
realised chain-level by the explicit map
$\PV^\bullet(X)\otimes\gl(V)\to\cE_\hCS(X,\gl(V))$ of
Theorem~\ref{thm:HKR-matrix-hCS} on local affine charts, globalised by
Calaque--Van den Bergh using $c_1(X)=0$.

(ii) \emph{The explicit presentation is the physics one.} $\cF^\hCS_X$ is
a concrete presentation of the abstract $E_3$ object: it comes with a
propagator, Feynman graphs, a BV Laplacian, one-loop anomaly computations,
and a compatible BCOV curving. The categorical presentation $\cF^{\HH}_X$
is abstractly unique but does not, by itself, carry a preferred propagator;
hCS provides the propagator and hence the Kontsevich associator point in
the $\Gtone$-torsor.

(iii) \emph{Curving aligns the presentations at the quantum level.} On a
compact CY$_3$ with $\int c_3(T_X)\neq 0$, quantum-level identification
of the two presentations requires the BCOV $B$-field (Theorem
\ref{thm:Atiyah-BCOV}) trivialising the $\kappa_{\mathrm{anom}}$-anomaly.
On $K3\times E$ the curving is inessential ($\kappa_{\mathrm{anom}}=0$ by
Wave-14 I1 Cycle 5).

(iv) \emph{The traces probe the common object.} $\DT$, $\GW$, $\PT$ of
Theorem~\ref{thm:DT-GW-MNOP} are $E_3$-factorisation-homology traces on
the common object $\Phi^{\mathrm{FA}}_3(X)$; they do not distinguish the
hCS and $\HH^\bullet(\cC)$ presentations --- they see only $\cF_X$. MNOP
and DT/PT are trace-identities, not functor-comparisons.
\end{corollary}

\begin{proof}[Proof]
(i) is the combination of Theorem~\ref{thm:KT-formality-d3} (KT formality
uniqueness up to $\Gtone$) and Theorem~\ref{thm:HKR-matrix-hCS} (HKR
$+$ matrix framing $=$ hCS at the chain level). (ii) is Wave-14 A2 Cycle
6: Costello--Li's BV propagator realises Kontsevich's configuration
integrals, pinning the Kontsevich associator. (iii) is Theorem
\ref{thm:Atiyah-BCOV}: Atiyah class and BCOV curving are the global
formality obstruction and its trivialisation respectively. (iv) is
Theorem~\ref{thm:DT-GW-MNOP}.
\end{proof}

\begin{remark}[The CG side-by-side]
\label{rmk:CG-side-by-side}
Costello--Francis--Gwilliam present the $E_n$ chain-level story through
factorisation algebras (Costello--Gwilliam~2017, Volumes I and II); Francis
presents the $(\infty,1)$-categorical story through higher Deligne and
the tangent complex of $E_n$-algebras (Francis~2013~\emph{Compos.\ Math.}
149; Lurie \emph{HA}~\S5). The two presentations of Corollary
\ref{cor:hcs-geometric-resolution} are exactly these two sides: hCS is
the Costello--Gwilliam chain-level factorisation-algebra presentation with
explicit propagator and Feynman graphs; $\HH^\bullet(\Perf(X))$ is the
Francis--Lurie $(\infty,1)$-categorical presentation with higher-Deligne
$E_3$-structure. CLAUDE.md's equal-status rule for chain-level and
$(\infty,1)$-categorical mathematics is here a theorem: the two
presentations are quasi-isomorphic $E_3$-objects by Theorem
\ref{thm:KT-formality-d3}, chain-level and $(\infty,1)$-categorical
access points to the same Stage-1 object $\Phi^{\mathrm{FA}}_3(X)$.
\end{remark}

\begin{remark}[Why this synthesis matters for the treatise]
\label{rmk:why-matters-treatise}
The three worked examples of \S\ref{wn:subsec:jordan-triple},
\S\ref{wn:subsec:resolved-conifold}, \S\ref{wn:subsec:K3xE} each exhibit
both presentations at their own level of rigour: $\ccC^3$ has both a hCS
realisation (Costello--Gwilliam Vol.~II \S10; Wave-14 J3) and a categorical
presentation via the Jordan triple quiver
(Theorem~\ref{thm:HKR-matrix-hCS}); the conifold has both a hCS
factorisation algebra and a quiver-with-potential category
(Szendroi~2008); $K3\times E$ has a categorical Davison BPS presentation
and an open hCS presentation (the BCOV anomaly computation of Wave-14 I1
supplies the compatibility). The Stage-1 object
$\Phi^{\mathrm{FA}}_3(X)$ exists uniformly across all three examples;
the chiral Yangian / $\cW_{1+\infty}$ output on the spectral line is the
Stage-2 specialisation.
\end{remark}

\subsubsection{Status summary}
\label{wn:subsubsec:hcs-vs-cat-status}

\begin{description}
\item[\ClaimStatusTheorem] Theorem~\ref{thm:KT-formality-d3}
(KT formality for $\cD_3$; Kontsevich~1997 $L_\infty$; Tamarkin~1998
$\Ger$-lift; Fresse/Lurie Koszul-duality upgrade to $E_3$;
Calaque--Van den Bergh~2007 CY globalisation with $c_1=0$).

\item[\ClaimStatusTheorem] Theorem~\ref{thm:HKR-matrix-hCS} (HKR $+$
matrix framing $=$ hCS on $\ccC^3$; Swan~1996; Yekutieli~2002 for HKR;
Costello--Li~2016 for the BV quantisation producing Kontsevich deformation
quantisation; Schiffmann--Vasserot~2013 for the CoHA = Maurer--Cartan
shadow).

\item[\ClaimStatusTheorem] Theorem~\ref{thm:Atiyah-BCOV} (Atiyah class
obstruction; Kapranov~1999; BCOV curving; Costello--Li~2012 chain-level
and $K3\times E$ K\"unneth vanishing of $c_3$; Wave-14 I1 for the
one-loop BV anomaly identification).

\item[\ClaimStatusTheorem] Theorem~\ref{thm:DT-GW-MNOP} parts on
$\DT$ (Thomas~1998; MNOP~2006), $\GW$ (Kontsevich~1995;
Behrend--Fantechi~1997), $\PT$ (Pandharipande--Thomas~2009), and the
MNOP conjecture proved for toric CY$_3$; Oberdieck--Pixton~2017 on
$K3\times E$ primitive class. \ClaimStatusConjectured\ MNOP in full
non-toric generality.

\item[\ClaimStatusScoped] Corollary~\ref{cor:hcs-geometric-resolution}
(synthesis): (i) and (ii) and (iv) are theorem under the scope of
Theorems~\ref{thm:KT-formality-d3}--\ref{thm:DT-GW-MNOP}; (iii) is
scoped to compact CY$_3$ with explicit BCOV trivialisation, open in
full generality for $\int c_3 \neq 0$.
\end{description}

\paragraph{Cross-references.} Wave-12 U1
(\texttt{wave12\_u1\_two\_stage\_functor.tex}) for the two-stage
factorisation $\Phi_3 = \Sp\circ\Phi^{\mathrm{FA}}_3$; Wave-14 A2
(\texttt{wave14\_a2\_KT\_formality\_kontsevich.tex}) for the
$\Gtone$-torsor vs.\ contractible-iso-classes distinction and the
Kontsevich-associator identification; Wave-14 I1
(\texttt{wave14\_i1\_hCS\_nonabelian\_anomaly.tex}) for the one-loop
BV anomaly computation on $\ccC^3$ and $K3\times E$; Wave-14 I2
(\texttt{wave14\_i2\_Linf\_Obs\_hCS.tex}) for the $L_\infty$-structure
beyond the minimal model and Atiyah-class obstruction on compact CY$_3$;
Wave-14 I3 (\texttt{wave14\_i3\_deformation\_hCS.tex}) for the
$E_3$-Hochschild deformation complex and holographic image; Wave-14 J3
(\texttt{wave14\_j3\_framed\_E3\_PROP\_action.tex}) for the framed
$E_3$-operad action via Fulton--MacPherson; Wave-14 J5
(\texttt{wave14\_j5\_chiral\_Koszul\_hCS.tex}) for the chiral Koszul
duality $U^\chir(\cE_\hCS)^! \simeq \CE^*_{\dbar,\chir}(\cE_\hCS, \cO)$.

\paragraph{Primary sources engaged.} Kontsevich~1997
\texttt{arXiv:q-alg/9709040}; Tamarkin~1998 \texttt{arXiv:math/9803025};
Drinfeld~1990 (quasi-Hopf, $\Gtone$); Willwacher~2014
\texttt{arXiv:1009.1654}; Fresse~2017 (\emph{Homotopy of Operads},
AMS SURV~217); Lurie \emph{Higher Algebra} \S5.2--5.5;
Swan~1996 (HKR); Yekutieli~2002 \texttt{arXiv:math/0007082};
Calaque--Van den Bergh~2007 \texttt{arXiv:0708.2725};
Kapranov~1999 \texttt{arXiv:math/9704009} (Atiyah class $=l_3^{\min}$);
Costello~2011 \emph{Renormalization and EFT};
Costello--Gwilliam~2017 \emph{Factorization Algebras in QFT}, Vols.~I--II;
Costello--Li~2012 \texttt{arXiv:1201.4522} (BCOV);
Costello--Li~2016 \texttt{arXiv:1606.00365} (hCS and Kontsevich
quantisation); BCOV~1994 \emph{Comm.\ Math.\ Phys.}~165;
Pantev--To\"en--Vaqui\'e--Vezzosi~2013 \texttt{arXiv:1111.3209}
(PTVV shifted symplectic); Kontsevich--Soibelman~2008
\texttt{arXiv:0811.2435} (CY$_3$ categories); Thomas~1998
\texttt{arXiv:math/9806111} (DT); Maulik--Nekrasov--Okounkov--Pandharipande~2006
\texttt{arXiv:math/0312059} (MNOP I); Pandharipande--Thomas~2009
\texttt{arXiv:0707.2348} (PT pairs); Oberdieck--Pixton~2017
\texttt{arXiv:1706.10100} (K3$\times$E reduced DT); Francis~2013
\texttt{arXiv:1110.5442} (higher Deligne); Ayala--Francis~2015
\texttt{arXiv:1212.5553} (non-abelian Poincar\'e duality).

%======================================================================
% End of Wave-14 B5/5 synthesis section.
%======================================================================
